\scanpage{206}
\mspage{122}
% archive sheet number: 205
% element: display | top, line 1 | x equals colimit x i
\[
x=\varinjlim_{i\in\unsure{I|\bar I}}x_i
\]
% element: text | top, lines 2-3 | les x i dans M pi
les $x_i$ dans $M_\pi$ ($\varinjlim$ \inserted{dans $M$} du foncteur
% element: diagram | upper third, line 4 | slice to M pi
\diagram{\begin{tikzcd}[ampersand replacement=\&]
I=M_\pi/x\arrow[r]\&M_\pi
\end{tikzcd}}
). \deleted{Donc} Quitte à remplacer $I$ par une catégorie ordonnée \inserted{$I'$} cofinale, cela prouve la partie qualitative de c. Supposons enfin $x\in\deleted{\unsure{\mathrm{Filt}^{\pi'}|\mathrm{Ob}}}\unsure{\mathrm{Ob}}\,M_{\pi'}$, où $\pi'>\pi$, \deleted{\unsure{on a donc}} montrons que l'on a alors
% element: deletion | interlinear above bound | M pi included in Fi
\deleted{$M_\pi\subset\unsure{\mathrm{Fi}|\mathrm{Filt}}$}
% element: display | middle, line 11 | cardinal bounds for I zero and I prime zero
\[
\operatorname{card}I_0,\operatorname{card}I'_{\unsure{0}}\leq\unsure{\pi'{}^{\pi^*}|\pi'{}^{\pi^\sharp}|\pi^{\pi^*}|\pi^{\pi'}|\pi'{}^\pi}
\tag*{(4.5.9.1)}
\]
% element: margin note | right of bound | où I zero I prime zero
\margin{où $I_0,I'_0$ sont les \unsure{sous-cat.|sous-catégories} pleines réduites de $I,I'$ telles que $I_0\subset I$, $I'_0\subset I'$ soient des $\approx$,}
% slanted
% element: margin note | left of bound, diagonally | On ne le prouve que
\margin{On ne le prouve que si $\pi\geq\pi_1$.}
% element: text | middle, lines 12-16 | ce qui \unsure{achèvera|achève} de prouver
ce qui \unsure{achèvera|achève} de prouver que l'on a une filtration cardinale. Il nous faut une majoration de cardinaux. \deleted{Il sera plus commode de \unsure{modifier|montrer} un peu}

% element: text | lower middle, line 17 | Lemme 4.5.10
\uline{Lemme} 4.5.10. Soit $\unsure{\pi^*|\pi^+|\pi^\sharp}$ un cardinal $\geq\pi_1$. Alors on a
% element: display | lower third, line 19 | cardinal Fl Filt skeleton bound
\[
\operatorname{card}\mathrm{Fl}\bigl((\mathrm{Filt}^{\pi^*})_0\bigr)
\leq\deleted{\unsure{\pi^{\pi_0}|\pi^*}}\,2^{\pi^*}.
\tag*{(4.5.10.1)}
\]
% slanted
% element: margin note | lower left, diagonally | Donner énoncé commun
\margin{\unsure{Donner énoncé} \unsure{commun pour} 4.5.11, et prouver d'abord 4.5.11.}
% element: text | foot, lines 20-25 | On a en effet
On a en effet, en désignant par $\unsure{\widehat C(\pi^*)|\widehat C(\leq\pi^*)}_{\deleted{\unsure{0}}}$ la sous-catégorie \inserted{str.} pleine de $\widehat C$ formée des $F:\unsure{C^\circ|C^0}\longrightarrow\mathrm{Ens}$ avec $\operatorname{card}F(x)\leq\pi^*$ pour \deleted{\unsure{tt}} $x$ dans $C$ (i.e. $F$ \unsure{se factorise|se factorisant} par $\mathrm{Ens}(\leq\pi^*)$), et par $\widehat C(\pi^*)_0$ une
